\subsection{Survival of every finite order}\label{sec:survival}

The defect law $\Gamma_t$ witnesses $Q(\eps,r)\in\Iexp_t$, which is Theorem~\ref{thm:membership}. The three feasible sets require successively more of $\Gamma_t$.

\subsubsection{The Navascués--Wolfe conditions}
Symmetry is Lemma~\ref{lem:symmetry}. For the diagonal law, let $R_\ell$ be the union of the three lines of the diagonal triangle $\Delta_{\ell\ell\ell}$:
\begin{equation}\label{eq:Rl}
R_\ell=\{(\ell,\ell,k):k\in[t]\}\cup\{(\ell,j,\ell):j\in[t]\}\cup\{(i,\ell,\ell):i\in[t]\}.
\end{equation}
Equivalently, $R_\ell$ is the set of cells with at least two coordinates equal to $\ell$; it has $3t-2$ cells.

\begin{lemma}[Diagonal supports are mutually disjoint]\label{lem:diag}
If $\ell\ne m$ then $R_\ell\cap R_m=\emptyset$. Consequently the $t$ diagonal triangles depend on mutually disjoint families of independent bits, are mutually independent under $\Gamma_t$, and their joint law is $Q(\eps,r)^{\otimes t}$.
\end{lemma}
\begin{proof}
A cell has three coordinates; it cannot have two of them equal to $\ell$ and two equal to $m\ne\ell$. The private bits of distinct diagonal triangles are distinct. Mutual independence of functions of disjoint families of independent variables follows, and Lemma~\ref{lem:triangle-law} identifies each factor.
\end{proof}

The lemma proves the tensor-power diagonal law, so $Q(\eps,r)\in\INW_t$.

\subsubsection{The ancestral-independence conditions}
Every injectable set is a subset of some copied triangle $\Delta_{ijk}$ (Definition~\ref{def:AI}), so its prescribed marginal is a projection of $Q(\eps,r)$ and is matched by Lemma~\ref{lem:triangle-law}. A union of pairwise ancestrally independent injectable sets is, by Lemma~\ref{lem:disjoint} applied to the union, a family of outputs depending on pairwise disjoint families of independent bits, hence mutually independent, with the prescribed product law. This proves $Q(\eps,r)\in\IAI_t$. In particular the higher-degree relations that Navascués and Wolfe mention but do not impose, such as $\Law(A^{12},B^{21},C^{12})=P_A\otimes P_B\otimes P_C$, hold for $\Gamma_t$: the three observations $A^{12},B^{21},C^{12}$ have ancestor sets $\{X_1,Z_2\}$, $\{X_2,Y_1\}$ and $\{Z_1,Y_2\}$. Their lines $(1,2,\cdot)$, $(2,\cdot,1)$ and $(\cdot,1,2)$ are pairwise disjoint.

\subsubsection{The recursively expressible conditions}
\begin{lemma}[Recursive prescriptions hold]\label{lem:expressible}
Every marginal prescribed by Definition~\ref{def:exp} equals the corresponding marginal of $\Gamma_t$.
\end{lemma}
\begin{proof}
The preceding argument proves the AI prescriptions. The root-sink lemma (Lemma~\ref{lem:rootsink}), applied to the triangle, gives every recursively expressible prescription.
\end{proof}

This proves $Q(\eps,r)\in\Iexp_t$ and completes the proof of Theorem~\ref{thm:membership}. Theorem~\ref{thm:main} follows by combining Theorem~\ref{thm:membership} at $r=r_t=(1-\eps_t)^{t-1}$ with Lemma~\ref{lem:violation}.
